\frontchapter{Abstract}

Drone-based radar search-and-rescue requires real-time spatial georeferencing: translating local radar detections into global GPS coordinates by fusing the drone's inertial state with radar range-Doppler measurements. This thesis presents the design and verification of a deterministic radar-inertial georeferencing coprocessor on the STM32H753 (Arm Cortex-M7), built entirely at the bare-metal register level and verified using Antmicro's open-source Renode functional simulator with hardware-in-the-loop testing using real MAVLink telemetry data.

The coprocessor sits between the drone's flight controller (Pixhawk running PX4/ArduPilot) and the radar SoC's digital backend. It ingests high-rate MAVLink attitude and position telemetry over USART1 using DMA with zero-copy circular buffering, hardware-timestampes incoming radar frames using a 32-bit free-running timer at 1\,$\mu$s resolution, runs a 6-state Extended Kalman Filter (EKF) fusing inertial and radar measurements using CMSIS-DSP matrix operations, converts local polar radar coordinates to global WGS84 (lat/lon/alt) via quaternion-based coordinate transforms, and streams georeferenced target metadata to the ground station over SPI DMA.

The firmware is developed across eleven progressive case studies, each isolating a specific pipeline stage and the defect it exposed: a build-system path resolution failure, an IRQn suffix typo, linker bare-metal C library dependencies, DMA-to-CPU L1 cache coherency hazards, struct-packing alignment shifts, hardware timer prescaler shadow loading, interrupt priority inversion, lock-free ring buffer memory ordering, EKF matrix buffer aliasing, coordinate transform quaternion rotation direction, output frame length mismatches, and MAVLink parser single-packet state management. Each defect is reproduced, root-caused using Renode's deterministic virtual-time execution and system-bus diagnostics, and corrected.

Hardware-in-the-loop verification is performed by injecting real MAVLink v2 packets (generated using pymavlink with correct CRCs) into the Renode simulation via memory-mapped backdoors, running the full pipeline end-to-end, and verifying georeferenced output coordinates. The Renode STM32H7 model's divergences from physical silicon are explicitly documented: DMAMUX1 is not modeled (worked around via direct buffer injection), timer input capture is not supported (timestamps injected via backdoor), and CRC hardware is simplified (CRC verification skipped in simulation builds). The firmware is defensively engineered for physical silicon correctness in all cases where the emulator diverges.

\vspace{0.8cm}
\noindent\textbf{Keywords:} Embedded Systems; ARM Cortex-M7; STM32H7; Renode; Hardware Emulation; Kalman Filtering; Georeferencing; MAVLink; Sensor Fusion; Deterministic Real-Time Systems; DMA; Cache Coherency.